\documentclass[12pt,a4paper]{article}
\usepackage[utf8]{inputenc}
\usepackage{amsmath, amssymb}
\usepackage{graphicx}
\usepackage{hyperref}
\usepackage{geometry}
\usepackage{booktabs}
\geometry{margin=1in}

\title{\textbf{Intelligent Crop Yield Prediction and Agentic Farm Advisory System}}
\author{Group Name/Members}
\date{\today}

\begin{document}

\maketitle

\begin{abstract}
This report details the design and implementation of an AI-driven agricultural analytics system. The system employs advanced machine learning techniques, including Cross-Validation and hyperparameter tuning, to predict crop yield from historical farm, soil, and weather data. Furthermore, an agentic AI assistant utilizing LangGraph and open-source Large Language Models (LLMs) was developed to generate structured, actionable agronomic advice based on the predicted yield.
\end{abstract}

\section{Introduction}
Crop yield prediction is crucial for agricultural planning and food security. In this project, we utilize supervised learning methods, specifically ensemble models such as Random Forest and Gradient Boosting, to predict continuous crop yield (hg/ha) based on inputs like average rainfall, temperature, pesticide usage, and crop type. We further augment the system by integrating an agentic workflow that translates these numerical predictions into structured farm management recommendations \cite{langgraph_docs}.

\section{Methodology}

\subsection{Data Preprocessing and Feature Engineering}
The dataset, sourced from Kaggle \cite{kaggle_dataset}, contains 28,242 records encompassing 101 countries and 10 major crops. We performed label encoding for categorical variables (Area, Item) and applied standard scaling for continuous features to ensure model convergence and robustness.

\subsection{Model Training and Hyperparameter Tuning}
We evaluated multiple models, focusing on regression rather than classification, as the target variable (crop yield) is continuous. We employed \textit{GridSearchCV} with 3-fold Cross-Validation to optimize hyperparameters, such as the number of estimators and max depth for the Random Forest Regressor \cite{scikit_learn}.

\subsection{Evaluation Metrics}
To gain a comprehensive understanding of model performance, we utilized the following metrics:
\begin{itemize}
    \item \textbf{Mean Absolute Error (MAE):} Measures the average magnitude of the errors.
    \begin{equation}
        \text{MAE} = \frac{1}{n} \sum_{i=1}^{n} |y_i - \hat{y}_i|
    \end{equation}
    
    \item \textbf{Root Mean Squared Error (RMSE):} Penalizes larger errors, derived from the MSE loss function.
    \begin{equation}
        \text{RMSE} = \sqrt{\frac{1}{n} \sum_{i=1}^{n} (y_i - \hat{y}_i)^2}
    \end{equation}
    
    \item \textbf{R-squared ($R^2$):} Represents the proportion of variance explained by the model.
\end{itemize}

\section{Results}
Our ensemble methods significantly outperformed the baseline linear models. The feature importance plots generated by the Random Forest model revealed that pesticide usage and average rainfall were primary drivers of yield variability.

\begin{figure}[h!]
    \centering
    % \includegraphics[width=0.7\textwidth]{assets/residual_plot.png}
    \caption{\textit{Placeholder for Scatter Plot of Actual vs. Predicted Yield and Residuals.} Note: A confusion matrix is not applicable here as we explicitly refined our approach to treat yield prediction as a continuous regression problem rather than a classification task.}
    \label{fig:residuals}
\end{figure}

\section{Agentic AI Farm Advisory Assistant}
The yield prediction system is extended using an agentic AI workflow implemented via LangGraph. The agent retrieves the predicted yield and farm parameters, analyzes potential risk factors, and outputs a structured advisory report including:
\begin{itemize}
    \item Crop and Field Summary
    \item Yield Prediction Interpretation
    \item Identified Risk Factors
    \item Recommended Farming Actions
\end{itemize}

\section{Limitations and Future Work}
While the system provides valuable insights, several limitations exist:
\begin{itemize}
    \item \textbf{Data Granularity:} The current dataset relies on country-level averages. Incorporating micro-climate and precise soil sensor data would improve prediction accuracy.
    \item \textbf{LLM Hallucinations:} Despite structured prompts, open-source LLMs may occasionally generate generic advice. Future work will involve implementing Retrieval-Augmented Generation (RAG) over specific, localized agricultural extension documents.
    \item \textbf{Temporal Dynamics:} Integrating Recurrent Neural Networks (RNNs) or LSTMs could better capture the temporal relationships in seasonal weather changes \cite{deep_learning_agri}.
\end{itemize}

\section{Conclusion}
The integration of optimized ensemble regression models with an LLM-based agentic workflow creates a comprehensive tool for both yield forecasting and actionable farm management. By addressing model generalizability through Cross-Validation and refining the user interface, the system serves as a robust prototype for intelligent agricultural analytics.

\begin{thebibliography}{9}

\bibitem{kaggle_dataset}
Patel, R. (2020). \textit{Crop Yield Prediction Dataset}. Kaggle. Available at: \url{https://www.kaggle.com/datasets/patelris/crop-yield-prediction-dataset}

\bibitem{scikit_learn}
Pedregosa, F. et al. (2011). \textit{Scikit-learn: Machine Learning in Python}. Journal of Machine Learning Research, 12, 2825-2830.

\bibitem{langgraph_docs}
LangChain AI. (2024). \textit{LangGraph Documentation}. Available at: \url{https://python.langchain.com/docs/langgraph}

\bibitem{deep_learning_agri}
Kamilaris, A., \& Prenafeta-Boldú, F. X. (2018). \textit{Deep learning in agriculture: A survey}. Computers and Electronics in Agriculture, 147, 70-90.

\bibitem{random_forest_agri}
Dahikar, V. S., \& Rode, S. V. (2014). \textit{Agricultural crop yield prediction using artificial neural network approach}. International Journal of Innovative Research in Electrical, Electronics, Instrumentation and Control Engineering, 2(1), 683-686.

\end{thebibliography}

\end{document}
